% Chapter 2: Implementation Plan and Architecture
% Assignment 2 - Communications and Cybersecurity

\section{Implementation Plan and Architecture}

\subsection{(a) Final Plan to Instantiate Services}

\subsubsection{Deployed Services}

We implemented six services fulfilling Assignment 2 requirements:

\begin{table}[h]
\centering
\begin{tabular}{|l|l|p{6cm}|}
\hline
\textbf{Service} & \textbf{Technology} & \textbf{Purpose} \\
\hline
PKI (ca01) & OpenSSL 3.0 & Certificate issuance, CRL management \\
\hline
Web Server (web01) & Node.js + Express & Public web service with 2FA \\
\hline
Database (db01) & PostgreSQL 15 & Data storage with mTLS \\
\hline
SSH Server (ssh01) & OpenSSH 9.x & Remote access with certificates \\
\hline
Firewall (fw01) & nftables & Traffic filtering, default DENY \\
\hline
Reverse Proxy (nginx01) & Nginx 1.24 & HTTPS termination \\
\hline
\end{tabular}
\caption{Implemented services and technologies}
\end{table}

\subsubsection{Deployment Steps}

The deployment follows a dependency-driven sequence:

\begin{enumerate}
    \item \textbf{Phase 1 - Network Setup}: Create DMZ (172.18.0.0/16) and Internal (172.19.0.0/16) networks with static IP assignments
    
    \item \textbf{Phase 2 - PKI Initialization}: Generate Root CA → Intermediate CA → Server certificates → Client certificates → SSH CA
    
    \item \textbf{Phase 3 - Service Deployment}: 
    \begin{verbatim}
    docker-compose up -d db01 web01 fw01 nginx01
    cd SSH01 && sudo ./setup-ssh-vm.sh
    \end{verbatim}
    
    \item \textbf{Phase 4 - Verification}: Test connectivity, certificate validation, firewall rules, and configuration integrity
\end{enumerate}

\subsubsection{Changes from Assignment 1}

Key modifications for Assignment 2:

\begin{enumerate}
    \item \textbf{Network Segmentation}: Added DMZ (172.18.x) and Internal (172.19.x) networks with static IPs for firewall rules
    
    \item \textbf{SSH Deployment}: Moved from Docker to Kali VM for realistic remote access testing
    
    \item \textbf{CA Isolation}: Firewall blocks all ca01 traffic (internal only, certificate issuance completed pre-deployment)
    
    \item \textbf{Auth \& AuthZ}: Added JWT authentication, fail2ban (2 failures=15min), RBAC (admin/read/readwrite)
    
    \item \textbf{Integrity}: SHA-256 checksums with automatic (fw01, web01) and manual (db01, ssh01, ca01) verification
    
    \item \textbf{2FA}: TOTP-based with Google Authenticator (enabled for new users, disabled for admin)
\end{enumerate}

\subsection{(b) Scenario Setup Diagram}

This section provides diagrams illustrating service deployment and network architecture.

\subsubsection{Architecture Diagram}

\begin{figure}[h]
\centering
\begin{verbatim}
                    INTERNET
                        |
                   [Port 443]
                        |
    ┌───────────────────┴────────────────────┐
    │         DMZ Network (172.18.0.0/16)    │
    │                                         │
    │   ┌──────────┐         ┌──────────┐   │
    │   │ nginx01  │────────>│  web01   │   │
    │   │.18.0.4   │  :3000  │.18.0.3   │   │
    │   │(HTTPS)   │         │(HTTPS)   │   │
    │   └──────────┘         └────┬─────┘   │
    │        │                    │          │
    │        └──────┬─────────────┘          │
    │               │                        │
    │          ┌────┴─────┐                  │
    │          │   fw01   │                  │
    │          │.18.0.2   │                  │
    │          │.19.0.4   │                  │
    │          │(nftables)│                  │
    │          └────┬─────┘                  │
    └───────────────┼────────────────────────┘
                    │
    ┌───────────────┴────────────────────────┐
    │   Internal Network (172.19.0.0/16)     │
    │                                         │
    │   ┌──────────┐         ┌──────────┐   │
    │   │  web01   │────────>│  db01    │   │
    │   │.19.0.5   │  :5432  │.19.0.3   │   │
    │   │(mTLS)    │  mTLS   │(PostgreSQL)│  │
    │   └──────────┘         └──────────┘   │
    │                                         │
    │        ┌──────────┐                    │
    │        │  ca01    │ (BLOCKED by FW)    │
    │        │.19.0.2   │                    │
    │        │(PKI)     │                    │
    │        └──────────┘                    │
    └─────────────────────────────────────────┘

                   ┌──────────┐
                   │ ssh01    │ (Kali VM)
                   │(SSH CA)  │ 192.168.x.x
                   │:22       │
                   └──────────┘
\end{verbatim}
\caption{Network architecture with DMZ and Internal segmentation}
\end{figure}

\textbf{Legend:}
\begin{itemize}
    \item Solid lines (──) indicate allowed network connections
    \item Arrow direction (─>) shows traffic flow
    \item Port numbers indicate service endpoints
    \item IP addresses show static assignments
\end{itemize}

\subsubsection{Infrastructure Description}

\paragraph{Network Segmentation}

The infrastructure implements a two-tier network segmentation model:

\begin{enumerate}
    \item \textbf{DMZ Network (172.18.0.0/16)}
    \begin{itemize}
        \item Purpose: Hosts publicly accessible services
        \item Gateway: 172.18.0.1
        \item Services: nginx01 (reverse proxy), web01 (dual-homed), fw01 (dual-homed)
        \item External Access: Port 443 exposed to internet
        \item Security: Firewall rules strictly control traffic flow
    \end{itemize}
    
    \item \textbf{Internal Network (172.19.0.0/16)}
    \begin{itemize}
        \item Purpose: Hosts backend services and sensitive data
        \item Gateway: 172.19.0.1
        \item Services: db01 (database), ca01 (PKI), web01 (dual-homed), fw01 (dual-homed)
        \item Isolation: No direct external access; CA completely blocked
        \item Security: mTLS required for all connections
    \end{itemize}
\end{enumerate}

\subsubsection{(c) Data Flows and Service Inventory}

\begin{table}[h]
\centering
\begin{tabular}{|l|l|l|l|p{4.5cm}|}
\hline
\textbf{Flow} & \textbf{Protocol} & \textbf{Port} & \textbf{Encryption} & \textbf{Purpose} \\
\hline
Client→nginx01 & HTTPS & 443 & TLS 1.3 & User web access \\
\hline
nginx01→web01 & HTTPS & 3000 & TLS 1.3 & Reverse proxy \\
\hline
web01→db01 & PostgreSQL & 5432 & TLS + mTLS & Database queries \\
\hline
Admin→ssh01 & SSH & 22 & Ed25519 & Remote access \\
\hline
ANY→ca01 & - & - & BLOCKED & CA isolation \\
\hline
\end{tabular}
\caption{Data flows with ports, protocols, and encryption}
\end{table}

\textbf{Network Segmentation:} DMZ (172.18.x): nginx01 (.4), web01 (.3), fw01 (.2); Internal (172.19.x): ca01 (.2), db01 (.3), fw01 (.4), web01 (.5). Firewall enforces default DENY with specific allow rules.

\subsubsection{Virtualization Model}

\textbf{Container Orchestration:} Docker 24.x+ with Docker Compose for service orchestration, custom bridge networks with static IPs, persistent volumes for database/certificates/logs.

\begin{table}[h]
\centering
\begin{tabular}{|l|l|l|}
\hline
\textbf{Service} & \textbf{Base Image} & \textbf{Rationale} \\
\hline
nginx01, web01, db01 & Alpine Linux 3.x & Minimal (5MB), security-focused \\
\hline
fw01, ca01 & Debian 12 & Full-featured (nftables, OpenSSL) \\
\hline
ssh01 & Kali Linux 2024.x & Pre-configured SSH server \\
\hline
\end{tabular}
\caption{Operating systems and deployment rationale}
\end{table}

\textbf{Security Boundaries:} Network isolation (custom Docker networks), filesystem isolation (minimal volume mounts, read-only where possible), process isolation (minimal privileges), certificate-based trust (no shared secrets).

\section{Policies and Cryptography}

\subsection{(d) Policy Implementation}

\begin{table}[h]
\centering
\begin{tabular}{|l|l|p{6.5cm}|}
\hline
\textbf{Service} & \textbf{Policy} & \textbf{Implementation} \\
\hline
Web Server & 2FA & TOTP with speakeasy, Google Authenticator \\
Web Server & Integrity & SHA-256 checksums (server.js, package.json) \\
Web Server & Encryption & HTTPS/TLS 1.3 via nginx reverse proxy \\
\hline
Firewall & Default DENY & nftables policy drop, specific allow rules \\
Firewall & Integrity & SHA-256 on nftables.conf, entrypoint.sh \\
\hline
Database & RBAC & JWT roles (admin=write, read=read only) \\
Database & Integrity & SHA-256 on postgresql.conf, pg\_hba.conf \\
Database & Encryption & PostgreSQL TLS + mTLS with client certs \\
\hline
SSH & Certificates & Ed25519 SSH CA, user certificates \\
SSH & Integrity & SHA-256 on sshd\_config, ca.pub \\
SSH & Sessions & MaxSessions 5 in sshd\_config \\
\hline
PKI & Secure Access & Directory permissions 700, restricted access \\
PKI & Integrity & SHA-256 on CA database and configs \\
PKI & Lifecycle & Issue, verify, revoke with CRL \\
\hline
SIEM & Fail2ban & 2 failures = 15min ban, in-memory tracking \\
SIEM & Integrity & SHA-256 on server.js \\
SIEM & Encryption & HTTPS/TLS 1.3 for admin access \\
\hline
\end{tabular}
\caption{Policy implementation per service (Assignment 2 Table 1)}
\end{table}

\subsection{(e) Cryptographic Mechanisms}

\begin{table}[h]
\centering
\begin{tabular}{|l|l|p{6.5cm}|}
\hline
\textbf{Mechanism} & \textbf{Algorithm} & \textbf{Usage} \\
\hline
TLS & TLS 1.3 & All HTTPS connections (client→nginx, nginx→web01) \\
\hline
Key Exchange & ECDHE & TLS handshake with perfect forward secrecy \\
\hline
Symmetric Encryption & AES-256-GCM & Bulk data encryption in TLS sessions \\
\hline
Hash (TLS) & SHA-384 & TLS 1.3 HMAC and cipher suites \\
\hline
Hash (Integrity) & SHA-256 & Configuration file checksums \\
\hline
Hash (Password) & bcrypt (cost=10) & User password storage in database \\
\hline
Certificates (TLS) & RSA 4096-bit & X.509 server/client certificates \\
\hline
Certificates (SSH) & Ed25519 & SSH user certificate authentication \\
\hline
2FA & TOTP (RFC 6238) & Time-based one-time passwords \\
\hline
JWT & HMAC-SHA256 & Token signing for session management \\
\hline
Database Auth & SCRAM-SHA-256 & PostgreSQL password authentication \\
\hline
\end{tabular}
\caption{Employed cryptographic mechanisms}
\end{table}

\section{Testing and Evaluation}

\subsection{(f) Test Plan and Results}

\subsubsection{Authentication Testing}

\textbf{Test 1: Admin login without 2FA}
\begin{verbatim}
curl -k -X POST https://localhost/login \
  -H "Content-Type: application/json" \
  -d '{"username":"admin","password":"adminpassword"}'
\end{verbatim}

\textbf{Result:} SUCCESS - Received JWT token with role='admin'
\begin{verbatim}
{"token":"eyJhbGciOiJIUzI1NiIsInR5cCI6IkpXVCJ9...",
 "user":{"username":"admin","role":"admin"}}
\end{verbatim}

\subsubsection{Authorization Testing (RBAC)}

\textbf{Test 2: Admin write access}
\begin{verbatim}
curl -k -X POST https://localhost/data \
  -H "Authorization: Bearer <admin-token>" \
  -H "Content-Type: application/json" \
  -d '{"title":"Test","content":"Admin post"}'
\end{verbatim}

\textbf{Result:} SUCCESS - Data created
\begin{verbatim}
{"id":5,"title":"Test","content":"Admin post","created_by":"admin"}
\end{verbatim}

\textbf{Test 3: User write access (should fail)}
\begin{verbatim}
curl -k -X POST https://localhost/data \
  -H "Authorization: Bearer <user-token>" \
  -H "Content-Type: application/json" \
  -d '{"title":"Test","content":"User post"}'
\end{verbatim}

\textbf{Result:} SUCCESS (403 as expected) - Permission denied
\begin{verbatim}
{"error":"Permission denied. READ-ONLY access."}
\end{verbatim}

\subsubsection{Fail2ban Testing}

\textbf{Test 4: Multiple failed login attempts}
\begin{verbatim}
# Attempt 1
curl -k -X POST https://localhost/login \
  -d '{"username":"admin","password":"wrong"}'
# Response: {"error":"Invalid credentials","remainingAttempts":2}

# Attempt 2
curl -k -X POST https://localhost/login \
  -d '{"username":"admin","password":"wrong"}'
# Response: {"error":"Invalid credentials","remainingAttempts":1}

# Attempt 3
curl -k -X POST https://localhost/login \
  -d '{"username":"admin","password":"wrong"}'
# Response: {"error":"Account locked after 2 failed attempts..."}
\end{verbatim}

\textbf{Result:} SUCCESS - Account locked for 15 minutes after 2 failed attempts

\subsubsection{Firewall Testing}

\textbf{Test 5: CA isolation}
\begin{verbatim}
docker exec fw01 nft list ruleset | grep -A 5 "CA-ISOLATION"
\end{verbatim}

\textbf{Result:} SUCCESS - Rules block all traffic to 172.19.0.2 (ca01)
\begin{verbatim}
ip daddr 172.19.0.2 log prefix "[FW-DENY-CA-ACCESS]" drop
\end{verbatim}

\subsubsection{SSH Certificate Testing}

\textbf{Test 6: SSH connection with certificate}
\begin{verbatim}
ssh -i PKI/ssh-ca/user-certs/admin-key admin@192.168.x.x
\end{verbatim}

\textbf{Result:} SUCCESS - Connected using certificate authentication (no password prompt)

\subsubsection{Configuration Integrity Testing}

\textbf{Test 7: Integrity verification}
\begin{verbatim}
docker exec web01 /usr/local/bin/config-integrity.sh verify
docker exec fw01 /usr/local/bin/integrity-check.sh verify
\end{verbatim}

\textbf{Result:} SUCCESS - All checksums matched baseline

\subsection{(g) Performance Measurement}

\textbf{Measurement Tool:} docker stats (resource utilization), curl -w (response time)

\subsubsection{Resource Utilization}

\begin{table}[h]
\centering
\begin{tabular}{|l|r|r|}
\hline
\textbf{Service} & \textbf{CPU \%} & \textbf{Memory (MB)} \\
\hline
nginx01 & 0.1\% & 8 MB \\
web01 & 0.5\% & 45 MB \\
db01 & 0.3\% & 25 MB \\
fw01 & 0.0\% & 3 MB \\
\hline
\textbf{Total} & \textbf{0.9\%} & \textbf{81 MB} \\
\hline
\end{tabular}
\caption{Resource usage under idle conditions}
\end{table}

\subsubsection{Response Time Analysis}

\begin{verbatim}
# Measure HTTPS response time
for i in {1..10}; do
  curl -k -w "%{time_total}\n" -o /dev/null -s https://localhost/health
done | awk '{sum+=$1} END {print "Average:",sum/NR*1000,"ms"}'
\end{verbatim}

\textbf{Results:}
\begin{itemize}
    \item Average response time: 18ms
    \item TLS handshake overhead: ~3ms (compared to HTTP baseline)
    \item Impact: +20\% latency for encryption
\end{itemize}

\subsubsection{Conclusion}

Security measures introduce minimal performance overhead. TLS 1.3 adds approximately 20\% latency but provides essential end-to-end encryption. Resource usage is minimal (81 MB total memory, <1\% CPU), making the infrastructure suitable for production deployment. The fail2ban mechanism operates in-memory with negligible performance impact.


